\begin{lemma}
\label{lemma-pseudo-coherent-hocolim}
Let $X$ be a quasi-compact and quasi-separated scheme.
Let $K \in D(\mathcal{O}_X)$. The following are equivalent
\begin{enumerate}
\item $K$ is pseudo-coherent, and
\item $K = \text{hocolim} K_n$ where
$K_n$ is perfect and $\tau_{\geq -n}K_n \to \tau_{\geq -n}K$
is an isomorphism for all $n$.
\end{enumerate}
\end{lemma}

\begin{proof}
The implication (2) $\Rightarrow$ (1) is true on any ringed space.
Namely, assume (2) holds. Recall that a perfect object of the derived
category is pseudo-coherent, see
Cohomology, Lemma \ref{cohomology-lemma-perfect}.
Then it follows from the definitions that
$\tau_{\geq -n}K_n$ is $(-n + 1)$-pseudo-coherent
and hence $\tau_{\geq -n}K$ is $(-n + 1)$-pseudo-coherent,
hence $K$ is $(-n + 1)$-pseudo-coherent. This is true for
all $n$, hence $K$ is pseudo-coherent, see
Cohomology, Definition \ref{cohomology-definition-pseudo-coherent}.

\medskip\noindent
Assume (1). We start by choosing an approximation
$K_1 \to K$ of $(X, K, -2)$ by a perfect complex $K_1$, see
Definitions \ref{definition-approximation-holds} and
\ref{definition-approximation} and
Theorem \ref{theorem-approximation}.
Suppose by induction we have
$$
K_1 \to K_2 \to \ldots \to K_n \to K
$$
with $K_i$ perfect such that
such that $\tau_{\geq -i}K_i \to \tau_{\geq -i}K$ is an isomorphism
for all $1 \leq i \leq n$. Then we pick $a \leq b$ as in
Lemma \ref{lemma-ext-from-perfect-into-bounded-QCoh}
for the perfect object $K_n$. Choose an approximation
$K_{n + 1} \to K$ of $(X, K, \min(a - 1, -n - 1))$.
Choose a distinguished triangle
$$
K_{n + 1} \to K \to C \to K_{n + 1}[1]
$$
Then we see that $C \in D_\QCoh(\mathcal{O}_X)$ has
$H^i(C) = 0$ for $i \geq a$. Thus by our choice of $a, b$
we see that $\Hom_{D(\mathcal{O}_X)}(K_n, C) = 0$.
Hence the composition $K_n \to K \to C$ is zero. Hence by
Derived Categories, Lemma \ref{derived-lemma-representable-homological}
we can factor $K_n \to K$ through $K_{n + 1}$
proving the induction step.

\medskip\noindent
We still have to prove that $K = \text{hocolim} K_n$.
This follows by an application of
Derived Categories, Lemma \ref{derived-lemma-cohomology-of-hocolim}
to the functors
$H^i( - ) : D(\mathcal{O}_X) \to \textit{Mod}(\mathcal{O}_X)$
and our choice of $K_n$.
\end{proof}
